\documentclass[11pt]{article}

\usepackage[margin=1in]{geometry}
\usepackage{amsmath,amssymb,amsthm}
\usepackage{mathtools}
\usepackage[T1]{fontenc}
\usepackage{lmodern}
\usepackage{enumitem}
\usepackage{microtype}

\setlength{\parindent}{0pt}
\setlength{\parskip}{5pt}

\newcommand{\cstir}[2]{c(#1,#2)}
\newcommand{\sstir}[2]{S(#1,#2)}

\begin{document}

\begin{center}
    {\Large\bfseries MATH 465 -- Homework 4 Solutions}\\[4pt]
    Due September 29, 2026
\end{center}

\textbf{Acknowledgment.} I consulted the homework prompt and no other books,
notes, or students. This write-up was prepared with generative AI assistance.

\textbf{Notation.} The signless Stirling number \(c(n,k)\) counts permutations
of \(\{1,\ldots,n\}\) having exactly \(k\) cycles. The Stirling number of the
second kind \(S(n,k)\) counts partitions of an \(n\)-element set into exactly
\(k\) nonempty unlabeled blocks.

\section*{1. A recurrence for \(c(n,k)\)}

Consider a permutation of \(\{1,\ldots,n\}\), and distinguish the element
\(n\). There are two disjoint cases.

\begin{itemize}[leftmargin=2em]
    \item If \(n\) is a one-element cycle, removing it leaves a permutation of
    \(\{1,\ldots,n-1\}\) with \(k-1\) cycles. There are \(c(n-1,k-1)\) such
    permutations.
    \item Otherwise, start with a permutation of \(\{1,\ldots,n-1\}\) with
    \(k\) cycles and insert \(n\) immediately after one of the existing
    \(n-1\) elements in its cycle notation. There are \(n-1\) choices of
    insertion position, and this operation does not change the number of
    cycles. Thus there are \((n-1)c(n-1,k)\) permutations in this case.
\end{itemize}

These two cases exhaust all possibilities, so
\[
    c(n,k)=c(n-1,k-1)+(n-1)c(n-1,k).
\]
The stated conventions make the same argument valid at the boundary values.

\section*{2. Three combinatorial identities}

\subsection*{(a) \(c(n,n-1)=\binom{n}{2}\)}

A permutation of \(n\) elements with \(n-1\) cycles has cycle deficit one:
relative to \(n\) fixed points, exactly one pair must be combined into a
2-cycle. Choose the two elements in \(\binom{n}{2}\) ways; their 2-cycle is
uniquely determined, and all other elements are fixed. Therefore
\[
    c(n,n-1)=\binom{n}{2}.
\]

\subsection*{(b) \(c(n,n-2)=2\binom{n}{3}+3\binom{n}{4}\)}

The cycle deficit is now two. There are exactly two possible cycle structures:

\begin{enumerate}[leftmargin=2em]
    \item One 3-cycle and all remaining elements fixed. Choose its three
    elements in \(\binom{n}{3}\) ways. On a specified set of three elements,
    there are \((3-1)!=2\) distinct 3-cycles, giving
    \(2\binom{n}{3}\) permutations.
    \item Two disjoint 2-cycles and all remaining elements fixed. Choose the
    four elements in \(\binom{n}{4}\) ways. A four-element set can be split
    into two unlabeled pairs in three ways, giving
    \(3\binom{n}{4}\) permutations.
\end{enumerate}

The cases are disjoint and exhaustive, so
\[
    c(n,n-2)=2\binom{n}{3}+3\binom{n}{4}.
\]

\subsection*{(c) \(S(n,n-2)=\binom{n}{3}+3\binom{n}{4}\)}

For a set partition into \(n-2\) blocks, the total block-size deficit from
the partition into singletons is two. Thus either there is one block of size
three, or there are two blocks of size two.

In the first case, choose the three elements in \(\binom{n}{3}\) ways. In the
second case, choose the four elements and partition them into two pairs. There
are three pairings of a specified four-element set, so this case contributes
\(3\binom{n}{4}\). Hence
\[
    S(n,n-2)=\binom{n}{3}+3\binom{n}{4}.
\]

\section*{3. Tourists at circular tables with drinks}

Let \(T_n\) be the number of seating-and-drink arrangements for \(n\)
distinct tourists. We use the auxiliary value \(T_0=1\), representing the
empty arrangement. Given an arrangement of \(n-1\) tourists, insert tourist
\(n\) in one of two ways:

\begin{itemize}[leftmargin=2em]
    \item Tourist \(n\) may form a new one-person table. The new table can
    choose any of the \(r\) drinks.
    \item Tourist \(n\) may join an existing table. To do this, choose the
    tourist who will be immediately to the right of tourist \(n\). There are
    \(n-1\) choices. Inserting tourist \(n\) in that position preserves the
    drink assigned to the table.
\end{itemize}

Thus every arrangement of \(n-1\) tourists produces exactly
\(r+(n-1)\) arrangements of \(n\) tourists. Conversely, deleting tourist
\(n\) recovers the unique preceding arrangement, whether \(n\) was alone or
not. Therefore
\[
    T_n=(r+n-1)T_{n-1}.
\]
Since \(T_0=1\), iterating gives
\[
    T_n=r(r+1)\cdots(r+n-1).
\]

\section*{4. A hockey-stick identity}

Consider all \((k+1)\)-element subsets of \(\{1,2,\ldots,n+1\}\). Classify
each subset according to its largest element. If its largest element is
\(j+1\), where \(k\leq j\leq n\), then the other \(k\) elements can be chosen
from \(\{1,\ldots,j\}\) in \(\binom{j}{k}\) ways. Hence the total number of
such subsets is
\[
    \sum_{j=k}^{n}\binom{j}{k}.
\]
On the other hand, choosing any \((k+1)\)-element subset directly gives
\(\binom{n+1}{k+1}\) possibilities. Since the two counts describe the same
set of subsets,
\[
    \boxed{\displaystyle \sum_{j=k}^{n}\binom{j}{k}=\binom{n+1}{k+1}}.
\]
No recurrence or algebraic manipulation of binomial coefficients is needed.

\section*{5. The sum of cubes}

For every nonnegative integer \(j\),
\[
    j^3=6\binom{j}{3}+6\binom{j}{2}+\binom{j}{1}.
\]
Summing and applying the identity from Problem 4 gives
\begin{align*}
    \sum_{j=0}^{n}j^3
    &=6\binom{n+1}{4}+6\binom{n+1}{3}+\binom{n+1}{2}\\
    &=\frac{(n+1)n(n-1)(n-2)}{4}
      +(n+1)n(n-1)+\frac{n(n+1)}{2}\\
    &=\boxed{\frac{n^2(n+1)^2}{4}}
      =\boxed{\left(\frac{n(n+1)}{2}\right)^2}.
\end{align*}

\section*{6. Linear recurrences}

\subsection*{(a)}

The recurrence has characteristic equation
\[
    \lambda^3-2\lambda^2-\lambda+2
    =(\lambda-2)(\lambda-1)(\lambda+1).
\]
The roots are distinct, so
\[
    a_n=A2^n+B+C(-1)^n.
\]
Using \(a_0=4\), \(a_1=8\), and \(a_2=10\) gives
\[
    A+B+C=4,\qquad 2A+B-C=8,\qquad 4A+B+C=10.
\]
Subtracting the first equation from the third gives \(A=2\); then
\(B+C=2\) and \(B-C=4\), so \(B=3\) and \(C=-1\). Therefore
\[
    \boxed{a_n=2^{n+1}+3-(-1)^n\qquad(n\geq 0).}
\]

\subsection*{(b)}

The characteristic equation is
\[
    \lambda^3-9\lambda^2+27\lambda-27=(\lambda-3)^3.
\]
Thus the general solution is
\[
    a_n=(A+Bn+Cn^2)3^n.
\]
The initial values yield
\[
    A=1,\qquad 3(A+B+C)=3,\qquad 9(A+2B+4C)=27.
\]
Hence \(B+C=0\) and \(2B+4C=2\), which give \(B=-1\) and \(C=1\). Thus
\[
    \boxed{a_n=(n^2-n+1)3^n\qquad(n\geq 0).}
\]

\section*{7. A recurrence for \(a_n=(2n-1)3^n+4(-1)^n\)}

The term \((2n-1)3^n\) is a linear combination of \(3^n\) and \(n3^n\),
so it corresponds to a root \(3\) of multiplicity two. The term \((-1)^n\)
corresponds to the root \(-1\). Therefore an annihilating characteristic
polynomial is
\[
    (\lambda-3)^2(\lambda+1)
    =\lambda^3-5\lambda^2+3\lambda+9.
\]
Consequently, a third-order homogeneous recurrence is
\[
    \boxed{a_n=5a_{n-1}-3a_{n-2}-9a_{n-3}\qquad(n\geq 3).}
\]

\section*{8. A Fibonacci sum}

Consider tilings of a board of length \(n\) by monomers of length one and
dominoes of length two. Let \(T_n\) be the number of such tilings. Looking at
the first tile gives
\[
    T_n=T_{n-1}+T_{n-2},\qquad T_0=1,\quad T_1=1.
\]
Therefore \(T_n=f_{n+1}\).

If a tiling has exactly \(k\) dominoes, then it has \(n-2k\) monomers and
there are \(n-k\) tiles in total. Choosing which \(k\) of those tile positions
are dominoes gives \(\binom{n-k}{k}\) tilings. The possible values are
\(0\leq k\leq\lfloor n/2\rfloor\). Summing over \(k\),
\[
    \boxed{f_{n+1}=\sum_{k=0}^{\lfloor n/2\rfloor}
    \binom{n-k}{k}}.
\]

\section*{9. The Fibonacci addition identity}

Again use tilings by monomers and dominoes. Consider the boundary between
positions \(m\) and \(m+1\) in a board of length \(m+n\). Partition all tilings
into two disjoint classes.

If no domino crosses the boundary, the first segment of length \(m\) and the
second segment of length \(n\) can be tiled independently. This gives
\[
    f_{m+1}f_{n+1}
\]
tilings. If a domino crosses the boundary, remove that domino. The remaining
segments have lengths \(m-1\) and \(n-1\), which can be tiled in
\[
    f_m f_n
\]
ways. Every tiling belongs to exactly one of these classes, so
\[
    \boxed{f_{m+n+1}=f_{m+1}f_{n+1}+f_mf_n}.
\]

\section*{10. Compositions into odd parts}

Let \(a_n\) count compositions of \(n\) whose parts are all odd.

\subsection*{(a) The Fibonacci generating function}

Let
\[
    F(x)=\sum_{n\geq 0}f_nx^n.
\]
Using \(f_0=0\), \(f_1=1\), and \(f_n=f_{n-1}+f_{n-2}\) for \(n\geq2\),
\begin{align*}
    F(x)
    &=x+\sum_{n\geq2}(f_{n-1}+f_{n-2})x^n\\
    &=x+xF(x)+x^2F(x).
\end{align*}
Therefore
\[
    \boxed{F(x)=\frac{x}{1-x-x^2}}.
\]

\subsection*{(b) The generating function for odd-part compositions}

The generating function for one positive odd part is
\[
    x+x^3+x^5+\cdots=\frac{x}{1-x^2}.
\]
A nonempty composition is a nonempty sequence of such parts. Consequently,
\begin{align*}
    A(x)&=\sum_{n\geq1}a_nx^n
    =\frac{\dfrac{x}{1-x^2}}{1-\dfrac{x}{1-x^2}}\\
    &=\boxed{\frac{x}{1-x-x^2}}.
\end{align*}
Comparing with part (a), the coefficient of \(x^n\) is the same for every
\(n\geq1\), so
\[
    \boxed{a_n=f_n\qquad(n\geq1).}
\]

\subsection*{(c) A bijection}

Write each odd part uniquely as \(2q+1\). Replace that part by a string
consisting of \(q\) copies of the part 2 followed by one copy of the part 1:
\[
    2q+1\longmapsto (\underbrace{2,\ldots,2}_{q\text{ copies}},1).
\]
Apply this replacement to every part of an odd-part composition of \(n\),
concatenate the resulting strings, and then delete the final 1. The sum of
the parts decreases by exactly one, so the result is a composition of
\(n-1\) with all parts equal to 1 or 2. For example,
\[
    (5,3,1)\longmapsto (2,2,1,2,1,1)\longmapsto(2,2,1,2,1,).
\]

For the inverse, start with a composition of \(n-1\) into 1s and 2s, append
a final 1, and split the sequence immediately after each 1. Each block has
the form \((2,\ldots,2,1)\), so it corresponds uniquely to an odd part. These
two operations are inverse to each other, proving the bijection. (For \(n=1\),
the empty composition of 0 is used.)

Finally, compositions of \(n-1\) into 1s and 2s are exactly tilings of a
board of length \(n-1\) by monomers and dominoes, and hence are counted by
\(f_n\). The bijection therefore independently proves
\[
    a_n=f_n.
\]

\end{document}
